\documentclass[12pt,a4paper]{article}
\PassOptionsToPackage{table}{xcolor}
% \documentclass{report}
\usepackage[utf8]{inputenc}
\usepackage{xspace}
\usepackage{url}
\usepackage{hyperref}
\usepackage{fancyhdr}
\usepackage{cite}
\usepackage{pgfgantt}
\usepackage{todonotes}
\usepackage[icelandic,UKenglish]{babel}
\usepackage[UKenglish]{datetime}
\usepackage[T1]{fontenc}
\usepackage{graphicx}
\usepackage[table]{xcolor}
\usepackage{enumitem}% http://ctan.org/pkg/enumitem
\usepackage[gen]{eurosym}
\usepackage{multicol}
\graphicspath{ {./images/} }
\usepackage[titletoc,title]{appendix}
\usepackage{pdfpages}
% \usepackage{array}
% \usepackage{pbox}
% \usepackage{hanging}
\usepackage{fancyhdr}
% \usepackage{caption}
% \usepackage{enumitem}
\usepackage{fancyvrb}
\usepackage{color}
\usepackage[utf8]{inputenc}
\usepackage{graphicx}
\usepackage{listings}
\usepackage{url}

\usepackage{xspace}

\usepackage{rotating}
\usepackage{wrapfig}

\usepackage{cite}
\usepackage{textcomp}

\usepackage[gen]{eurosym}
\usepackage{hyperref}

\usepackage{todonotes}

\usepackage{gensymb}

% \usepackage{draftwatermark}
% \SetWatermarkText{draft v2}
% \SetWatermarkScale{6}

\fancypagestyle{plain}{ %
  \fancyhf{} % remove everything
  \renewcommand{\headrulewidth}{0pt} % remove lines as well
  \renewcommand{\footrulewidth}{0pt}
}

\hyphenation{EISCAT}

\newenvironment{MYitemize}{%
    %% \renewcommand{\labelitemi}{$\rightarrow$}%
    %% \renewcommand{\labelitemii}{$\circ$}%
    %% \renewcommand{\labelitemii}{$\rightarrow$}%
    %% \renewcommand{\labelitemiii}{$\rightarrow$}%
    %% \renewcommand{\labelitemiii}{\colblack $\cdot$}%
    \begin{itemize}}{\end{itemize}}

\newcommand{\bitm}{\begin{MYitemize}}
\newcommand{\eitm}{\end{MYitemize}}

\newcommand{\N}{NeIC\xspace}
\newcommand{\NP}{NeIC~E3DDS\xspace} % NeIC project

\newcommand{\ED}{EISCAT\_3D\xspace}
\newcommand{\EC}{EISCAT\xspace}
\newcommand{\ESA}{EISCAT Scientific Association\xspace}
\newcommand{\SA}{sub-array\xspace}
\newcommand{\SAs}{sub-arrays\xspace}
\newcommand{\RS}{radar site\xspace}
\newcommand{\RSs}{radar sites\xspace}
\newcommand{\OS}{on-site\xspace}
\newcommand{\OC}{operations centre\xspace}
\newcommand{\CC}{control centre\xspace}
\newcommand{\DC}{data centre\xspace}
\newcommand{\DCa}{dCache\xspace}
\newcommand{\DCs}{data centres\xspace}
\newcommand{\UAF}{User Analysis Facility\xspace}

\newcommand{\RB}{ring buffer\xspace}
\newcommand{\fsru}{first stage receive unit\xspace}
\newcommand{\FW}{file writer\xspace}
\newcommand{\SBF}{second beam former\xspace}
\newcommand{\NNB}{$2 \times 100$\xspace} %number of narrow beams

\newcommand{\NBW}{5~MHz\xspace} % narrow bandwidth
\newcommand{\WBW}{30~MHz\xspace} % wide bandwidth

\newcommand{\ramf}{Ramfjordmoen\xspace}
\newcommand{\controllatency}{less than $5$\xspace}
\newcommand{\CP}{C/C++\xspace}
\newcommand{\DI}{DIRAC\xspace}
\newcommand{\DRF}{DigitalRF\xspace}
\newcommand{\HDF}{HDF5\xspace}
\newcommand{\IB}{Infiniband\xspace}
\newcommand{\IBtwo}{200\,Gb/s~\IB}
\newcommand{\Ru}{Rucio\xspace}
\newcommand{\einfra}{e-infrastructure\xspace}
\newcommand{\neinfra}{national e-infrastructure\xspace}
\newcommand{\WR}{White Rabbit\xspace}
\newcommand{\UAP}{User access portal\xspace}

\newcommand{\bfmust}{\xspace{\bf must}\xspace}
\newcommand{\bfshould}{\xspace{\bf should}\xspace}
\newcommand{\Bfmust}{\xspace{\bf Must}\xspace}
\newcommand{\Bfshould}{\xspace{\bf Should}\xspace}
\newcommand{\bfshall}{\xspace{\bf shall}\xspace}
\newcommand{\Bfshall}{\xspace{\bf Shall}\xspace}

\newcommand{\FirstBF}{\xspace{\bf Shall}\xspace}

\newcommand{\visroot}{isualiz}
\newcommand{\vis}{\visroot{ation}\xspace}
\newcommand{\vise}{\visroot{e}\xspace}
% \newcommand{\Vis}{{Visualization}\xspace}
\newcommand{\threed}{{3-dimensional}\xspace}
\newcommand{\GD}{{GUISDAP}\xspace}

\newcommand{\gps}{{Gbit/s}\xspace}

\newcommand{\np}{national provider\xspace}
\newcommand{\nps}{\np{s}\xspace}

\newcommand\EatDot[1]{}

\newcommand{\pc}{prompt computing\xspace}

%\newcommand{\NBW}{5~MHz} % narrow bandwidth
%\newcommand{\WBW}{30~MHz} % wide bandwidth


\newcommand{\nnt}{NeIC NT1\xspace}

% \setlength{\topskip}{0mm}
\setlength{\headheight}{15pt}
% \setlength{\topmargin}{-5.4mm}
% \setlength{\textheight}{230mm}
\setlength{\textwidth}{180mm}
\setlength{\oddsidemargin}{-5.0mm}
% \setlength{\evensidemargin}{10.0mm}
% \setlength{\captionmargin}{7mm}

\title{
{\bf Deliverable Document 4} \\
The synergies with the \nnt operations}
\author{E3DDS Team~\footnote{
Anders Tjulin (EISCAT) {\tt anders.tjulin@eiscat.se};
Ari Lukkarinen (CSC) {\tt ari.lukkarinen@csc.fi};
Assar Westman (EISCAT) {\tt Assar.Westman@eiscat.se};
Carl-Fredrik Enell (EISCAT) {\tt carl-fredrik.enell@eiscat.se};
Dan Johan Jonsson (UiT) {\tt dan.jonsson@uit.no};
Janos Nagy (NSC) {\tt fconagy@nsc.liu.se};
Harri Hellgren (EISCAT) {\tt harri.hellgren@eiscat.se};
Ingemar H\"{a}ggstr\"{o}m (EISCAT) {\tt ingemar.haggstrom@eiscat.se};
Mattias Wadenstein (UmU) {\tt maswan@hpc2n.umu.se};
% Roy Dragseth (UiT) {\tt roy.dragseth@uit.no};
John White (NeIC) {\tt john.white@cern.ch}}}

\date{\today}

\begin{document}

\pagestyle{fancy}
\lhead{\bf E3DDS project}
\rhead{\bf 4: NT1 Synergies}

\maketitle
\par\noindent
\begin{minipage}{0.5\textwidth}
  \includegraphics[scale=0.18]{NEIC_logo_screen_black.pdf}
  %\vspace{-0.09in}
\end{minipage}
\begin{minipage}{0.5\textwidth}
  \hfill
  %\includegraphics[scale=0.25]{EISCAT3Dlogo1.pdf}
  % New official logo with green text
  \includegraphics[width=0.75\linewidth]{e3d-logo-green-500px}
\end{minipage}

\newpage
\tableofcontents
\newpage

\section{Executive Summary}
\label{exec-summ}

%\todo[inline]{Executive summary to be written at the end of the document writing process? Moved the current exec summary to introduction as that's what it was more like}

Deliverable 4: The synergies with the \nnt operations. Once decisions on the software to be deployed and run as a simulation of the \ED operations, the possibilities of inclusion to the NT1 operations are studied.The object to be delivered is a document that details the services for \ED operations and those that may or may not be managed by \nnt. This includes the services and service management (e.g. FitSM) required and the division of responsibilities.  Included are estimates of the staffing required to run these services and SLAs needed. 

The study of the synergies between \ED online and NT1 operations may start once the software decisions are made in Deliverable 1 in the sixth month of the project.

The process that follows is continuous and requires input from both \ED and NT1.

Document delivered. See section 3.2.

\section{Purpose}
\label{purpose}

\subsection{Intended Audience}

The intended audience of this document is primarily NeIC and the \ED project management and staff in order to understand 

\section{Introduction}
\label{intro}

\newlength{\mycolwidth}
\setlength{\mycolwidth}{3.0cm}
\begin{center}
\begin{table}
\rowcolors{2}{gray!25}{white}
\begin{tabular}{p{\mycolwidth} p{\mycolwidth} l l l l p{\mycolwidth}}
{\bf \tiny Service}  & {\bf \tiny Location} & \bf {\tiny Real-time} & {\bf \tiny Run by} & \bf {\tiny Owned by} & {\bf \tiny Priority} & {\bf \tiny Comments} \\
% 
%  \bf Service  & \bf Location & \bf Real-time & \bf Run by & \bf Owned by & \bf Priority & \bf Comments
\bf \tiny Service  & \bf \tiny Location & \bf \tiny Real-time & \bf \tiny Run by & \bf \tiny Owned by & \bf \tiny Priority & \bf \tiny Comments
\\
 \tiny Radar controller & \tiny Skibotn or DC & \tiny Yes & \tiny EISCAT & \tiny EISCAT & \tiny 1 High & \tiny Radar controller  system
\\
 \tiny Ring buffer + Beam former + File Writer + Fast Imaging & \tiny Skibotn or all sites & \tiny Yes & \tiny EISCAT & \tiny EISCAT & \tiny 1 High & \tiny RAM and fast disks
\\
 \tiny Prompt computing   & \tiny TBD Skibotn or DC & \tiny Yes & \tiny EISCAT & \tiny EISCAT & \tiny 1 High & \tiny Lag profiling (Level 1 to 2), parameter fitting (Level 2 to 3), visualisation (Level 3 on geographic coordinates)
\\
 \tiny Secondary data products & \tiny Co-located with prompt computing & \tiny No & \tiny Other & \tiny EISCAT & \tiny 2 Medium & \tiny Standard data products that don't need to be produced in realtime e.g. full imaging
\\
 \tiny Identity management & \tiny One master site. Other secondary sites. & \tiny No & \tiny Other & \tiny EISCAT & \tiny 2 Medium & \tiny Directory services for user identity management
\\
 \tiny Site data buffer  & \tiny Co-located with prompt computing & \tiny Yes & \tiny Other & \tiny EISCAT & \tiny 1 High & \tiny Fast disk storage for processing of data, levels 1b, 2, 3. Partition of long-term storage?
\\
 \tiny Cluster management etc & \tiny Sites & \tiny No & \tiny Other & \tiny EISCAT & \tiny 1 High & \tiny Configuration and software deployment on RBBF, Prompt computing, admin login nodes etc. 
\\
 \tiny User portal web GUI & \tiny TBD & \tiny No & \tiny Other & \tiny Other & \tiny 2 Medium & \tiny User data discovery and job submission portal. DIRAC prototype exists.
\\
 \tiny User job management & \tiny TBD & \tiny No & \tiny Other & \tiny Other & \tiny 2 Medium& \tiny Submit analysis jobs. DIRAC considered.
\\
 \tiny Data staging & \tiny Initially Skibotn storage & \tiny No & \tiny Other & \tiny Other & \tiny 2 Medium & \tiny Prototype uses DIRAC storage element and file transfer 
\\
 \tiny User compute & \tiny Initially using Prompt computing & \tiny No & \tiny Other & \tiny EISCAT & \tiny 2 Medium & \tiny Computing resources for user analysis portal.
\\
 \tiny User software deployment & \tiny Complements User analysis portal & \tiny No & \tiny Other & \tiny Other & \tiny 2 Medium & \tiny e.g. the DIRAC input sandbox
\\
 \tiny User software deployment & \tiny Complements User analysis portal & \tiny No & \tiny Other & \tiny Other & \tiny 3 Low & \tiny e.g. Notebooks Juypter or other? 
\\
 \tiny Data archive & \tiny Data Centres & \tiny No & \tiny Other & \tiny EISCAT & \tiny 2 Medium & \tiny Long-term data storage. Does EISCAT want to run a data archive?
\\
 \tiny Internal metadata catalog & \tiny Data Centres & \tiny Yes & \tiny Other & \tiny EISCAT & \tiny 1 High & \tiny Metadata master catalogue e.g. metadata from Level 1b files. Realtime.
\\
 \tiny User metadata catalog & \tiny Data Centres & \tiny No & \tiny Other & \tiny EISCAT & \tiny 2 Medium & \tiny Metadata for users. e.g. from higher level data products.
\\
 \tiny Data product repository & \tiny Close to long-term storage & \tiny No & \tiny Other & \tiny EISCAT & \tiny 2 Medium & \tiny 
\\
 \tiny Data product publishing & \tiny Close to long-term storage & \tiny No & \tiny Other & \tiny EISCAT & \tiny 3 Low & \tiny e.g. FAIR compliant Repository for published datasets. E.g. B2SHARE.
\\
 \tiny Monitoring & \tiny Sites & \tiny Yes & \tiny Other & \tiny EISCAT & \tiny 1 High & \tiny Environment and operational status. Possibly NAV software from Uninett.
\\ 
 \tiny Inventory & \tiny Skibotn & \tiny No & \tiny EISCAT & \tiny EISCAT & \tiny 1 High & \tiny List of all parts (preferably marked with QR codes) and updates to status.
\\
 \tiny Public IP routing & \tiny Sites & \tiny Yes & \tiny NORDUnet  & \tiny NORDUnet & \tiny 3 Low & \tiny Internet access for technical staff and guest instruments.
\\

% \end{tabular}
% \caption{The \einfra services required for \ED. \label{tab:services}}
% \end{table}
% \end{center}


%  \bf Service  & \bf Location & \bf Real-time & \bf Run by & \bf Owned by & \bf Priority & \bf Comments
\bf \tiny Service  & \bf \tiny Location & \bf \tiny Real-time & \bf \tiny Run by & \bf \tiny Owned by & \bf \tiny Priority & \bf \tiny Comments
\\
 \tiny Radar controller & \tiny Skibotn or DC & \tiny Yes & \tiny EISCAT & \tiny EISCAT & \tiny 1 High & \tiny Radar controller  system
\\
 \tiny Ring buffer + Beam former + File Writer + Fast Imaging & \tiny Skibotn or all sites & \tiny Yes & \tiny EISCAT & \tiny EISCAT & \tiny 1 High & \tiny RAM and fast disks
\\
 \tiny Prompt computing   & \tiny TBD Skibotn or DC & \tiny Yes & \tiny EISCAT & \tiny EISCAT & \tiny 1 High & \tiny Lag profiling (Level 1 to 2), parameter fitting (Level 2 to 3), visualisation (Level 3 on geographic coordinates)
\\
 \tiny Secondary data products & \tiny Co-located with prompt computing & \tiny No & \tiny Other & \tiny EISCAT & \tiny 2 Medium & \tiny Standard data products that don't need to be produced in realtime e.g. full imaging
\\
 \tiny Identity management & \tiny One master site. Other secondary sites. & \tiny No & \tiny Other & \tiny EISCAT & \tiny 2 Medium & \tiny Directory services for user identity management
\\
 \tiny Site data buffer  & \tiny Co-located with prompt computing & \tiny Yes & \tiny Other & \tiny EISCAT & \tiny 1 High & \tiny Fast disk storage for processing of data, levels 1b, 2, 3. Partition of long-term storage?
\\
 \tiny Cluster management etc & \tiny Sites & \tiny No & \tiny Other & \tiny EISCAT & \tiny 1 High & \tiny Configuration and software deployment on RBBF, Prompt computing, admin login nodes etc. 
\\
 \tiny User portal web GUI & \tiny TBD & \tiny No & \tiny Other & \tiny Other & \tiny 2 Medium & \tiny User data discovery and job submission portal. DIRAC prototype exists.
\\
 \tiny User job management & \tiny TBD & \tiny No & \tiny Other & \tiny Other & \tiny 2 Medium& \tiny Submit analysis jobs. DIRAC considered.
\\
 \tiny Data staging & \tiny Initially Skibotn storage & \tiny No & \tiny Other & \tiny Other & \tiny 2 Medium & \tiny Prototype uses DIRAC storage element and file transfer 
\\
 \tiny User compute & \tiny Initially using Prompt computing & \tiny No & \tiny Other & \tiny EISCAT & \tiny 2 Medium & \tiny Computing resources for user analysis portal.
\\
 \tiny User software deployment & \tiny Complements User analysis portal & \tiny No & \tiny Other & \tiny Other & \tiny 2 Medium & \tiny e.g. the DIRAC input sandbox
\\
 \tiny User software deployment & \tiny Complements User analysis portal & \tiny No & \tiny Other & \tiny Other & \tiny 3 Low & \tiny e.g. Notebooks Juypter or other? 
\\
 \tiny Data archive & \tiny Data Centres & \tiny No & \tiny Other & \tiny EISCAT & \tiny 2 Medium & \tiny Long-term data storage. Does EISCAT want to run a data archive?
\\
 \tiny Internal metadata catalog & \tiny Data Centres & \tiny Yes & \tiny Other & \tiny EISCAT & \tiny 1 High & \tiny Metadata master catalogue e.g. metadata from Level 1b files. Realtime.
\\
 \tiny User metadata catalog & \tiny Data Centres & \tiny No & \tiny Other & \tiny EISCAT & \tiny 2 Medium & \tiny Metadata for users. e.g. from higher level data products.
\\
 \tiny Data product repository & \tiny Close to long-term storage & \tiny No & \tiny Other & \tiny EISCAT & \tiny 2 Medium & \tiny 
\\
 \tiny Data product publishing & \tiny Close to long-term storage & \tiny No & \tiny Other & \tiny EISCAT & \tiny 3 Low & \tiny e.g. FAIR compliant Repository for published datasets. E.g. B2SHARE.
\\
 \tiny Monitoring & \tiny Sites & \tiny Yes & \tiny Other & \tiny EISCAT & \tiny 1 High & \tiny Environment and operational status. Possibly NAV software from Uninett.
\\ 
 \tiny Inventory & \tiny Skibotn & \tiny No & \tiny EISCAT & \tiny EISCAT & \tiny 1 High & \tiny List of all parts (preferably marked with QR codes) and updates to status.
\\
 \tiny Public IP routing & \tiny Sites & \tiny Yes & \tiny NORDUnet  & \tiny NORDUnet & \tiny 3 Low & \tiny Internet access for technical staff and guest instruments.
\\

\end{tabular}
\caption{The \einfra services required for \ED. \label{tab:services}}
\end{table}
\end{center}

\section{\ED Operations}

The services shown in Table~\ref{tab:services} are those estimated, at the time of writing, to be necessary for \ED to: acquire data; process online and offline; produce data products in order to steer the radar; serve data to users for analysis.

The radar data is used to control subsequent radar transmission and receiver look directions.
Therefore the online computing services must run in real-time while the radar is transmitting.

A service that is noted as ``run by" \EC is a service that required specialized in-house knowledge that is unlikely to be found elsewhere.
The ``owned by" designation for a service refers to the management and ownership of the data in the service.
For example, the ``Site data buffer" disk array could be operated by a non-\EC entity whereas the choices on the data to retain or remove from the disks are made by \EC.
Another example could be the ``User portal web GUI" where typically operating a complex web portal is better left to experienced experts and the \EC users would access as clients where data retention etc is handled by central portal policies.

The priority of a service is defined by: {\bf 1 High:} Without this service the radar system CANNOT be run; {\bf 2 Medium:} The radar can still run but data cannot be got out of the system; {\bf 3 Low:} Radar can run and data extracted but analysis and publishing not possible.

The following scenarios for running the \ED \einfra are considered.
In these scenarios, the standard assumptions are:
\bitm
\item That the \ED data is replicated to two redundant data archives.
\item The \ED radars are planned to run 24/7~\footnote{this is a big caveat...}
\eitm

In each scenario the following factors are considered.
The {\bf Risks} and {\bf Benefits} overall in having the various organizations running the services.
The {\bf Cost Efficiency} refers to how much effort an organization will have to expend to run a set of service(s).
{\bf Reliability} describes whether the service up-time can be maintained in order to process and serve the data from the expected duty cycle of the radar operations.
In these scenarios the ``users" are the consumers of the \ED data and the comments on {\bf User Satisfaction} evaluates how this would be affected by each scenario.

\subsection{Scenario 1}

In this scenario \EC runs all the computing on-line and off-line services noted in Table~\ref{tab:services} except for the second redundant external data archive i.e. the copy of the \ED data.

\subsubsection*{Risks}
\bitm
\item Large amount of FTEs required to run all services.
\item Recruiting lots of skilled IT staff into \ESA might fail.
\item Estimate of effort? 
\eitm

\subsubsection*{Benefits}
\bitm
\item Hardware and staff are controlled through one set of \EC processes.
\eitm

\subsubsection*{Cost Efficiency}
\bitm
\item It may not be cost efficient for one organization to run all the services. Finding people with enough skills to run and maintain a wide variety of services usually difficult.
\item Replication of already existing services run by other \einfra projects e.g. EGI Checkin~\cite{egi-checkin}.
\eitm

\subsubsection*{Reliability}
\bitm
\item Reliability might suffer if enough people are not able to be hired to cover for holidays/sickness/job changes.
\eitm

\subsubsection*{User Satisfaction}
\bitm
\item User satisfaction may be high as users are able to contact directly to \EC personnel for help with problems or questions.
\item User satisfaction might be negatively impacted by a learning period for how to run many new services.
\eitm

%%%%%%%%%%%%%%%%%%%%%%%%%%%%%%%%%%%%%%%%%%%%%%%%%%%%%%%%%%%%%%%%%%%%%%%%%%%%%

\subsection{Scenario 2}
\label{ssec:scen2}

The \nnt programme runs all computing services listed in Table~\ref{tab:services} except the \ED inventory system and FSRUs under a service contract.

\subsubsection*{Risks}
\bitm
\item Large amount of FTEs required to run all services.
\item Recruiting lots of skilled IT staff into NT1 might fail.
\eitm

\subsubsection*{Benefits}
\bitm
\item \EC controls the services most familiar for them i.e. the FSRUs and the inventory databases.
\item Hardware and staff for running all other services are controlled through one set of \nnt processes.
\item Recruiting into the \nnt will be easier.
\eitm

\subsubsection*{Cost Efficiency}
\bitm
\item As with Scenario 1, it may not be cost efficient for one organization (\nnt) to run all the services. Finding people with enough skills to run and maintain a wide variety of services usually difficult.
\item Also runs the risk of replication of already existing services run by other \einfra projects.
\item Some services are not suited to be run by NT1.
\eitm

\subsubsection*{Reliability}
\bitm
\item The reliability of some the services should be improved as the \ED addition would be an increment on already existing \nnt services.
\item Services that are not part of the \nnt core services would not necessarily be more reliable in this scenario.
\item \EC would be free to concentrate on running their critical services.
\eitm

% Risks and benefits roughly the same.
% Recruiting easier.
% Cost benefits the same.
% Some services not suited to be run by NT1. Replicating services.
% Boundaries would have to be very tightly defined. SLAs.
% RBBF nodes... also tricky to run. Custom build of nodes. Tuning.
% Internal metadata catalogue?
% Experiment database... boundaries? Tricky boundary conditions... SLAs. Agents that program the FSRUs will pull data. 

\subsubsection*{User Satisfaction}
\bitm
\item \ED User satisfaction may be lowered as users will be required to contact an external organization for help with problems or questions. Not directly to other \EC colleagues.
\item User satisfaction might be negatively impacted by a learning period for how to run many new services.
\item After the initial learning period, the user satisfaction should improve to a level higher than if \EC was to run all services (Scenario 1).
\eitm

%%%%%%%%%%%%%%%%%%%%%%%%%%%%%%%%%%%%%%%%%%%%%%%%%%%%%%%%%%%%%%%%%%%%%%%%%

\subsection{Scenario 3}
\label{ssec:scen3}

In this scenario \EC EISCAT runs the high-priority services and outsources medium/low to their partners.
Table~\ref{tab:services:s3} shows the services that \EC would run.
\begin{center}
\begin{table}
\rowcolors{2}{gray!25}{white}
\begin{tabular}{p{\mycolwidth} p{\mycolwidth} l l l l p{\mycolwidth}}
{\bf \tiny Service}  & {\bf \tiny Location} & \bf {\tiny Real-time} & {\bf \tiny Run by} & \bf {\tiny Owned by} & {\bf \tiny Priority} & {\bf \tiny Comments} \\
% 
%  \bf Service  & \bf Location & \bf Real-time & \bf Run by & \bf Owned by & \bf Priority & \bf Comments
\bf \tiny Service  & \bf \tiny Location & \bf \tiny Real-time & \bf \tiny Run by & \bf \tiny Owned by & \bf \tiny Priority & \bf \tiny Comments
\\
 \tiny Radar controller & \tiny Skibotn or DC & \tiny Yes & \tiny EISCAT & \tiny EISCAT & \tiny 1 High & \tiny Radar controller  system
\\
 \tiny Ring buffer + Beam former + File Writer + Fast Imaging & \tiny Skibotn or all sites & \tiny Yes & \tiny EISCAT & \tiny EISCAT & \tiny 1 High & \tiny RAM and fast disks
\\
 \tiny Prompt computing   & \tiny TBD Skibotn or DC & \tiny Yes & \tiny EISCAT & \tiny EISCAT & \tiny 1 High & \tiny Lag profiling (Level 1 to 2), parameter fitting (Level 2 to 3), visualisation (Level 3 on geographic coordinates)
\\
 \tiny Secondary data products & \tiny Co-located with prompt computing & \tiny No & \tiny Other & \tiny EISCAT & \tiny 2 Medium & \tiny Standard data products that don't need to be produced in realtime e.g. full imaging
\\
 \tiny Identity management & \tiny One master site. Other secondary sites. & \tiny No & \tiny Other & \tiny EISCAT & \tiny 2 Medium & \tiny Directory services for user identity management
\\
 \tiny Site data buffer  & \tiny Co-located with prompt computing & \tiny Yes & \tiny Other & \tiny EISCAT & \tiny 1 High & \tiny Fast disk storage for processing of data, levels 1b, 2, 3. Partition of long-term storage?
\\
 \tiny Cluster management etc & \tiny Sites & \tiny No & \tiny Other & \tiny EISCAT & \tiny 1 High & \tiny Configuration and software deployment on RBBF, Prompt computing, admin login nodes etc. 
\\
 \tiny User portal web GUI & \tiny TBD & \tiny No & \tiny Other & \tiny Other & \tiny 2 Medium & \tiny User data discovery and job submission portal. DIRAC prototype exists.
\\
 \tiny User job management & \tiny TBD & \tiny No & \tiny Other & \tiny Other & \tiny 2 Medium& \tiny Submit analysis jobs. DIRAC considered.
\\
 \tiny Data staging & \tiny Initially Skibotn storage & \tiny No & \tiny Other & \tiny Other & \tiny 2 Medium & \tiny Prototype uses DIRAC storage element and file transfer 
\\
 \tiny User compute & \tiny Initially using Prompt computing & \tiny No & \tiny Other & \tiny EISCAT & \tiny 2 Medium & \tiny Computing resources for user analysis portal.
\\
 \tiny User software deployment & \tiny Complements User analysis portal & \tiny No & \tiny Other & \tiny Other & \tiny 2 Medium & \tiny e.g. the DIRAC input sandbox
\\
 \tiny User software deployment & \tiny Complements User analysis portal & \tiny No & \tiny Other & \tiny Other & \tiny 3 Low & \tiny e.g. Notebooks Juypter or other? 
\\
 \tiny Data archive & \tiny Data Centres & \tiny No & \tiny Other & \tiny EISCAT & \tiny 2 Medium & \tiny Long-term data storage. Does EISCAT want to run a data archive?
\\
 \tiny Internal metadata catalog & \tiny Data Centres & \tiny Yes & \tiny Other & \tiny EISCAT & \tiny 1 High & \tiny Metadata master catalogue e.g. metadata from Level 1b files. Realtime.
\\
 \tiny User metadata catalog & \tiny Data Centres & \tiny No & \tiny Other & \tiny EISCAT & \tiny 2 Medium & \tiny Metadata for users. e.g. from higher level data products.
\\
 \tiny Data product repository & \tiny Close to long-term storage & \tiny No & \tiny Other & \tiny EISCAT & \tiny 2 Medium & \tiny 
\\
 \tiny Data product publishing & \tiny Close to long-term storage & \tiny No & \tiny Other & \tiny EISCAT & \tiny 3 Low & \tiny e.g. FAIR compliant Repository for published datasets. E.g. B2SHARE.
\\
 \tiny Monitoring & \tiny Sites & \tiny Yes & \tiny Other & \tiny EISCAT & \tiny 1 High & \tiny Environment and operational status. Possibly NAV software from Uninett.
\\ 
 \tiny Inventory & \tiny Skibotn & \tiny No & \tiny EISCAT & \tiny EISCAT & \tiny 1 High & \tiny List of all parts (preferably marked with QR codes) and updates to status.
\\
 \tiny Public IP routing & \tiny Sites & \tiny Yes & \tiny NORDUnet  & \tiny NORDUnet & \tiny 3 Low & \tiny Internet access for technical staff and guest instruments.
\\

% \end{tabular}
% \caption{The \einfra services required for \ED. \label{tab:services}}
% \end{table}
% \end{center}

 \tiny Radar controller &
 \tiny Skibotn or DC &
 \tiny Yes &
 \tiny EISCAT &
 \tiny EISCAT &
 \tiny 1 High &
 \tiny Radar controller system \\
 \tiny Ring buffer + Beam former + File Writer + Fast Imaging & \tiny Skibotn or all sites & \tiny Yes & \tiny EISCAT & \tiny EISCAT & \tiny 1 High & \tiny RAM and fast disks \\
 \tiny Prompt computing   & \tiny TBD Skibotn or DC & \tiny Yes & \tiny EISCAT & \tiny EISCAT & \tiny 1 High & \tiny Lag profiling (Level 1 to 2), parameter fitting (Level 2 to 3), visualisation (Level 3 on geographic coordinates) \\
 \tiny Site data buffer  & \tiny Co-located with prompt computing & \tiny Yes & \tiny Other & \tiny EISCAT & \tiny 1 High & \tiny Fast disk storage for processing of data, levels 1b, 2, 3. Partition of long-term storage? \\
 \tiny Cluster management etc & \tiny Sites & \tiny No & \tiny Other & \tiny EISCAT & \tiny 1 High & \tiny Configuration and software deployment on RBBF, Prompt computing, admin login nodes etc.  \\
 \tiny Internal metadata catalog & \tiny Data Centres & \tiny Yes & \tiny Other & \tiny EISCAT & \tiny 1 High & \tiny Metadata master catalogue e.g. metadata from Level 1b files. Realtime. \\
 \tiny Monitoring & \tiny Sites & \tiny Yes & \tiny Other & \tiny EISCAT & \tiny 1 High & \tiny Environment and operational status. Possibly NAV software from Uninett. \\
 \tiny Inventory & \tiny Skibotn & \tiny No & \tiny EISCAT & \tiny EISCAT & \tiny 1 High & \tiny List of all parts (preferably marked with QR codes) and updates to status. \\

\end{tabular}
\caption{The High-priority \einfra services run by \EC in Scenario 3. \label{tab:services:s3}}
\end{table}
\end{center}
Table~\ref{tab:services:s3a} gives the low and medium priority services run by \EC partners.
It can be seen that these services are those that \EC is most likely to have expertise and interest to run at 
(Dirac at STFC, data staging and data archive at NT1, batch compute at HPC2N, metadata catalog at CSC, software management at UiT, etc).

\begin{center}
\begin{table}
\rowcolors{2}{gray!25}{white}
\begin{tabular}{p{\mycolwidth} p{\mycolwidth} l l l l p{\mycolwidth}}
{\bf \tiny Service}  & {\bf \tiny Location} & \bf {\tiny Real-time} & {\bf \tiny Run by} & \bf {\tiny Owned by} & {\bf \tiny Priority} & {\bf \tiny Comments} \\
% 
%  \bf Service  & \bf Location & \bf Real-time & \bf Run by & \bf Owned by & \bf Priority & \bf Comments
\bf \tiny Service  & \bf \tiny Location & \bf \tiny Real-time & \bf \tiny Run by & \bf \tiny Owned by & \bf \tiny Priority & \bf \tiny Comments
\\
 \tiny Radar controller & \tiny Skibotn or DC & \tiny Yes & \tiny EISCAT & \tiny EISCAT & \tiny 1 High & \tiny Radar controller  system
\\
 \tiny Ring buffer + Beam former + File Writer + Fast Imaging & \tiny Skibotn or all sites & \tiny Yes & \tiny EISCAT & \tiny EISCAT & \tiny 1 High & \tiny RAM and fast disks
\\
 \tiny Prompt computing   & \tiny TBD Skibotn or DC & \tiny Yes & \tiny EISCAT & \tiny EISCAT & \tiny 1 High & \tiny Lag profiling (Level 1 to 2), parameter fitting (Level 2 to 3), visualisation (Level 3 on geographic coordinates)
\\
 \tiny Secondary data products & \tiny Co-located with prompt computing & \tiny No & \tiny Other & \tiny EISCAT & \tiny 2 Medium & \tiny Standard data products that don't need to be produced in realtime e.g. full imaging
\\
 \tiny Identity management & \tiny One master site. Other secondary sites. & \tiny No & \tiny Other & \tiny EISCAT & \tiny 2 Medium & \tiny Directory services for user identity management
\\
 \tiny Site data buffer  & \tiny Co-located with prompt computing & \tiny Yes & \tiny Other & \tiny EISCAT & \tiny 1 High & \tiny Fast disk storage for processing of data, levels 1b, 2, 3. Partition of long-term storage?
\\
 \tiny Cluster management etc & \tiny Sites & \tiny No & \tiny Other & \tiny EISCAT & \tiny 1 High & \tiny Configuration and software deployment on RBBF, Prompt computing, admin login nodes etc. 
\\
 \tiny User portal web GUI & \tiny TBD & \tiny No & \tiny Other & \tiny Other & \tiny 2 Medium & \tiny User data discovery and job submission portal. DIRAC prototype exists.
\\
 \tiny User job management & \tiny TBD & \tiny No & \tiny Other & \tiny Other & \tiny 2 Medium& \tiny Submit analysis jobs. DIRAC considered.
\\
 \tiny Data staging & \tiny Initially Skibotn storage & \tiny No & \tiny Other & \tiny Other & \tiny 2 Medium & \tiny Prototype uses DIRAC storage element and file transfer 
\\
 \tiny User compute & \tiny Initially using Prompt computing & \tiny No & \tiny Other & \tiny EISCAT & \tiny 2 Medium & \tiny Computing resources for user analysis portal.
\\
 \tiny User software deployment & \tiny Complements User analysis portal & \tiny No & \tiny Other & \tiny Other & \tiny 2 Medium & \tiny e.g. the DIRAC input sandbox
\\
 \tiny User software deployment & \tiny Complements User analysis portal & \tiny No & \tiny Other & \tiny Other & \tiny 3 Low & \tiny e.g. Notebooks Juypter or other? 
\\
 \tiny Data archive & \tiny Data Centres & \tiny No & \tiny Other & \tiny EISCAT & \tiny 2 Medium & \tiny Long-term data storage. Does EISCAT want to run a data archive?
\\
 \tiny Internal metadata catalog & \tiny Data Centres & \tiny Yes & \tiny Other & \tiny EISCAT & \tiny 1 High & \tiny Metadata master catalogue e.g. metadata from Level 1b files. Realtime.
\\
 \tiny User metadata catalog & \tiny Data Centres & \tiny No & \tiny Other & \tiny EISCAT & \tiny 2 Medium & \tiny Metadata for users. e.g. from higher level data products.
\\
 \tiny Data product repository & \tiny Close to long-term storage & \tiny No & \tiny Other & \tiny EISCAT & \tiny 2 Medium & \tiny 
\\
 \tiny Data product publishing & \tiny Close to long-term storage & \tiny No & \tiny Other & \tiny EISCAT & \tiny 3 Low & \tiny e.g. FAIR compliant Repository for published datasets. E.g. B2SHARE.
\\
 \tiny Monitoring & \tiny Sites & \tiny Yes & \tiny Other & \tiny EISCAT & \tiny 1 High & \tiny Environment and operational status. Possibly NAV software from Uninett.
\\ 
 \tiny Inventory & \tiny Skibotn & \tiny No & \tiny EISCAT & \tiny EISCAT & \tiny 1 High & \tiny List of all parts (preferably marked with QR codes) and updates to status.
\\
 \tiny Public IP routing & \tiny Sites & \tiny Yes & \tiny NORDUnet  & \tiny NORDUnet & \tiny 3 Low & \tiny Internet access for technical staff and guest instruments.
\\

% \end{tabular}
% \caption{The \einfra services required for \ED. \label{tab:services}}
% \end{table}
% \end{center}

 \tiny Secondary data products & \tiny Co-located with prompt computing & \tiny No & \tiny Other & \tiny EISCAT & \tiny 2 Medium & \tiny Standard barch processing. \\
 \tiny Identity management & \tiny One master site. Other secondary sites. & \tiny No & \tiny Other & \tiny EISCAT & \tiny 2 Medium & \tiny e.g. EGI Checkin run by EGI partner? \\
 \tiny User portal web GUI & \tiny TBD & \tiny No & \tiny Other & \tiny Other & \tiny 2 Medium & \tiny DIRAC portal at RAL? \\
 \tiny User job management & \tiny TBD & \tiny No & \tiny Other & \tiny Other & \tiny 2 Medium& \tiny DIRAC portal at RAL? \\
 \tiny Data staging & \tiny Initially Skibotn storage & \tiny No & \tiny Other & \tiny Other & \tiny 2 Medium & \tiny NeIC NT1. \\
 \tiny User compute & \tiny Initially using Prompt computing & \tiny No & \tiny Other & \tiny EISCAT & \tiny 2 Medium & \tiny Academic Cloud provider. CSC? \\
 \tiny User software deployment & \tiny Complements User analysis portal & \tiny No & \tiny Other & \tiny Other & \tiny 2 Medium & \tiny DIRAC portal at RAL \\
 \tiny Data archive & \tiny Data Centres & \tiny No & \tiny Other & \tiny EISCAT & \tiny 2 Medium & \tiny NeIC NT1. \\
 \tiny User metadata catalog & \tiny Data Centres & \tiny No & \tiny Other & \tiny EISCAT & \tiny 2 Medium & \tiny NeIC NT1 \\
 \tiny Data product repository & \tiny Close to long-term storage & \tiny No & \tiny Other & \tiny EISCAT & \tiny 2 Medium & \tiny NeIC NT1 \\
 \tiny User software deployment & \tiny Complements User analysis portal & \tiny No & \tiny Other & \tiny Other & \tiny 3 Low & \tiny Partner University running Juypter or other?  \\
 \tiny Data product publishing & \tiny Close to long-term storage & \tiny No & \tiny Other & \tiny EISCAT & \tiny 3 Low & \tiny B2SHARE at CSC \\
 \tiny Public IP routing & \tiny Sites & \tiny Yes & \tiny NORDUnet  & \tiny NORDUnet & \tiny 3 Low & \tiny NORDUnet. \\

\end{tabular}
\caption{The Low and Medium-priority \einfra services run by partners in Scenario 3. \label{tab:services:s3a}}
\end{table}
\end{center}


\subsubsection*{Risks}
\bitm
\item Services running in different administrative regimes.
\eitm

\subsubsection*{Benefits}
\bitm
\item Services are run by organizations best suited and already operating such services.
\eitm

\subsubsection*{Cost Efficiency}
\bitm
\item Services are run by individuals already familiar with the services.
\eitm

\subsubsection*{Reliability}
\bitm
\item Services are run by individuals already familiar with the services.
\item Organizations can take in the \ED load as a fractional increase on their existing workload.
\eitm

\subsubsection*{User Satisfaction}
\bitm
\item May be impacted as users may have to contact service providers outside the \EC organization.
  \item Hopefully the service operators will be able to solve problems quickly due to experience.
\eitm

\subsection{Scenario 4}
\label{ssec:scen4}

EISCAT runs most, but outsources things with obvious synergies at other places (data archive, batch computing, web portals).
Other places... NT1 and other partners.

User Portal (with someone already running DIRAC for other communities)

Identity management (EGI/INFN/etc that handles federated identities)

User Compute (with someone that runs large batch computing)

User software deployment (with someone that runs scientific computing)

Data Archive, Data Staging, and Site Data Buffer could be synergistically run as one storage (NT1)

Data product repository and publishing

File catalog (with someone that is already running Rucio)

\subsubsection*{Risks}
\bitm
\item 
\eitm

\subsubsection*{Benefits}
\bitm
\item 
\eitm

\subsubsection*{Cost Efficiency}
\bitm
\item 
\eitm

\subsubsection*{Reliability}
\bitm
\item 
\eitm

\subsubsection*{User Satisfaction}
\bitm
\item 
\eitm
\subsection{Scenario 5}
\label{ssec:scen5}

Like Section~\ref{ssec:scen3} but NT1 operates EISCAT's inhouse resources under a service contract.

How is this different to 2? ... EISCAT outsources low and medium to partners. High is operated by NT1 on under a service contract.

\subsubsection*{Risks}
\bitm
\item 
\eitm

\subsubsection*{Benefits}
\bitm
\item 
\eitm

\subsubsection*{Cost Efficiency}
\bitm
\item 
\eitm

\subsubsection*{Reliability}
\bitm
\item 
\eitm

\subsubsection*{User Satisfaction}
\bitm
\item 
\eitm

\subsection{Scenario 6}
\label{ssec:scen6}

Like Section~\ref{ssec:scen4} but NT1 operates EISCAT's inhouse resources under a service contract.

NT1 would operate the services that do not have obvious synergy with the EISCAT ``partners".

\subsubsection*{Risks}
\bitm
\item 
\eitm

\subsubsection*{Benefits}
\bitm
\item 
\eitm

\subsubsection*{Cost Efficiency}
\bitm
\item 
\eitm

\subsubsection*{Reliability}
\bitm
\item 
\eitm

\subsubsection*{User Satisfaction}
\bitm
\item 
\eitm
% end of the document...
\newpage
\bibliography{main}{}
\bibliographystyle{unsrt}


\end{document}
